\section{Revisiting the noise distribution}
\label{sec:distribution}
Even though the resulting public key and ciphertext sizes at $k=7$ remain relatively small, 
increasing $k$ further is not possible without affecting correctness. Instead, one could reduce the failure probability at large $k$ by increasing the modulus $q$.
We note that all variants of \textsc{Kyber} fix $q=3329$, which is explicitly justified in the submission documentation to NIST. The modulus
$q=3329$ is the small prime satisfying the NTT-enabling condition $n \mid (q-1)$, notably weaker than the $2n \mid (q-1)$ one might expect, 
which \textsc{Kyber} deliberately does not require~\cite{kyberround2}. It was reduced from the round-1 value $7681$ to recover the
bandwidth spent by dropping public-key compression and co-tuned to balance ciphertext size, failure probability, and security. 

Since ML-KEM's standardization, highly optimized and high-assurance implementations have emerged across software and hardware platforms \cite{mlkem-native,fvk,SP:AOPPSS25}.
Changing $q$ would mean not being able to reuse any existing optimized routines in software or hardware, which sacrifices an important engineering advantage 
and results in a different scheme altogether.

This is not to say that $k$ is the only parameter that can be modified. We focus on changes that preserve the \emph{essence} of Kyber:
the ring $\mathbb{Z}_q[X]/(X^{256}+1)$, the NTT implementation and its
constants, and the existing analysis framework. The module
rank $k$ qualifies trivially. So does the error distribution. 
Since the noise only affects the error term and does not interact with the arithmetic structure or the NTT, 
it can be replaced without altering the core implementation. What it changes are only the MLWE hardness estimate and the 
decryption-failure probability, both of which can be recomputed through the same scripts, over the same ring. 


In this sense, the error distribution is another degree of freedom that moves both security and corectness without changing what the scheme is at its core. 
This is the direction we pursue and where our main contribution lies.

To scale past $k=7$, we need a noise knob that pushes the failure-driving part of the error back down as $k$ grows, while preserving 
\textsc{Kyber}'s \emph{essence}. Equally important, it should retain the cheap, branchless, constant-time sampling that motivates \textsc{Kyber}'s use of a 
centered binomial in the first place. This requirement rules out simply using a discrete Gaussian, as in FrodoKEM.

\subsection{A parameterized noise distribution}
We introduce a one-parameter family $f(t)$ on $\{-2,-1,0,1,2\}$, the same
support as $\mathrm{CBD}_2$, defined by
\[
  \Pr[\pm 2] = 2^{-t}, \qquad
  \Pr[\pm 1] = \tfrac14, \qquad
  \Pr[0] = \tfrac12 - 2^{-(t-1)}.
\]
The error at $\pm 1$ is fixed. The single parameter $t$ acts only on the
$\pm 2$ tail, thinning it geometrically. A sample is drawn from uniform
bits $a_1,\dots,a_t$ as
\[
  x = (2a_t - 1)\Bigl(1 - a_1 + 2\textstyle\prod_{i=1}^{t-1} a_i\Bigr),
\]
which is branchless and therefore samples in constant time; we
machine-check both constant-time (CT) and speculative constant-time (SCT)
in Jasmin~\cite{CCS:ABBBGL17} (see Section~\ref{sec:implementation}). The standard deviation
$\sigma(t) = \sqrt{0.5 + 2^{3-t}}$ decreases from $1$ at $t=4$ toward
$\sqrt{0.5}$ as $t$ grows.

Crucially, $f(t)$ generalizes Kyber's own noise distribution. At $t=4$, the probabilities become
$$
  \Pr[\pm 2] = \frac{1}{16}, \qquad
  \Pr[\pm 1] = \frac14, \qquad
  \Pr[0] = \frac38,
$$
exactly matching $\mathrm{CBD}_2$, the noise distribution used by ML-KEM-768 and -1024. As $t \to \infty$, the $\pm 2$ mass vanishes and $f(t)$ approaches $\mathrm{CBD}_1$.

\subsection{Why thinning the tail is the efficient trade}
A key property of $f(t)$ is that increasing $t$ reduces the probability of the $\pm2$ terms exponentially while leaving the $\pm1$ mass unchanged. As
defined above, we have
$$
\Pr_{f(t)}[\pm2] = 2^{-t},
\qquad
\Pr_{f(t)}[\pm1] = \frac14.
$$
So, $t$ is the tunable parameter for thinning the tail of
the distribution without changing the bulk $\pm1$ terms. This gives finer control over the noise than the standard centered binomial
choices. In particular, $\mathrm{CBD}_1$ and $\mathrm{CBD}_2$ have standard deviations $\sqrt{0.5}$ and $1$, 
respectively, with no centered binomial parameter providing an intermediate distribution. As shown above, 
$f(t)$ ranges from $\mathrm{CBD}_2$ toward $\mathrm{CBD}_1$ as $t$ increases, giving us finer control over the noise.

This finer control is useful as the module rank $k$ increases. Increasing
$k$ increases the amount of error accumulated during encryption and
therefore makes the decryption-failure constraint more restrictive. Rather
than moving directly from $\mathrm{CBD}_2$ to $\mathrm{CBD}_1$, we can
increase $t$ as $k$ grows, reducing the $\pm2$ contribution
while keeping the $\pm1$ contribution fixed.

\subsection{Selecting $t$}
Since $t$ is the number of uniform bits per sample, vectorization-friendly
values are natural, especially considering the AVX2 implementations. As shown above, $t=4$ recovers the $\mathrm{CBD}_2$ sampler, so it serves
as our anchor. Notice that $t=8$ consumes exactly one PRF byte per coefficient, meaning it
admits one-byte-per-coefficient SIMD packing without cross-byte bit-splitting. This makes $t=8$ a very convenient choice.
Our third choice, $t=6$, sits in the middle and does not share the byte-alignment of $t=4$ or $t=8$. Sampling four
coefficients consumes $6 \times 4 = 24$ bits, i.e.\ three bytes of PRF output, which is not as natural to implement
as the byte-aligned $t=4$ and $t=8$ cases. It is nonetheless interesting relative to $t=8$, for two reasons.
First, it retains more of the $\pm2$ tail, giving a higher standard deviation and more security per dimension.
Second, it spends fewer PRF bytes per coefficient ($0.75$ vs.\ $1$), which is relevant on platforms \emph{without} Keccak acceleration.
For each choice of $t$, we pick the highest $k$ satisfying the failure probability bound of $2^{-128}$, staying within the communication budget of
FrodoKEM640. 

We use FrodoKEM-640 as our baseline for parameter selection and comparison throughout the paper because it is the 
smallest standardized FrodoKEM parameter set and therefore imposes the tightest communication budget of 9,616 public-key 
bytes and 9,700 ciphertext bytes. This gives us 
a restrictive baseline for asking how much security margin \textsc{Kaiburr} can achieve within the communication cost 
of a conservative lattice-based KEM.

The three parameter sets we carry forward are $f(4)$, $f(6)$, and $f(8)$, at $k=7$, $18$, and $24$
respectively, which follow from these choices
(Table~\ref{tab:kaiburr-params}).

\begin{table}[H]
  \centering
  \caption{\textsc{Kaiburr} parameter sets. Classical/quantum security (Core-SVP) and failure probability computed with the Kyber team's scripts; other parameters as in ML-KEM, compression removed.}
  \label{tab:kaiburr-params}
  \begin{tabular}{lccccccc}
    \toprule
    Scheme & $k$ & noise & failure & class. & quant. & pk (B) & ct (B) \\
    \midrule
    kaiburr-1792 & 7  & $f(4)$ & $2^{-133.7}$ & 479  & 434  & 2720 & 3072 \\
    kaiburr-4608 & 18 & $f(6)$ & $2^{-134.1}$ & 1281 & 1162 & 6944 & 7296 \\
    kaiburr-6144 & 24 & $f(8)$ & $2^{-139.9}$ & 1717 & 1557 & 9248 & 9600 \\
    \bottomrule
  \end{tabular}
\end{table}

\subsection{The ceiling of the family}
The family $f(t)$ is not unbounded. The lightest noise it can
produce is its limit $f(\infty)$, with $\Pr[\pm 1]=\tfrac14$,
$\Pr[0]=\tfrac12$, $\Pr[\pm 2]=0$, and minimal variance
$\sigma^2=0.5$. At $n=256$ and $q=3329$, the failure probability crosses
$2^{-128}$ at $k=30$, so $k=29$ is the largest rank that $f(t)$ can
support (Table~\ref{tab:theoretical-ceiling}). Reaching higher
security levels would require another change, such as modifying $q$, $n$, or the noise family itself, 
moving beyond the \textsc{Kyber} structure we seek to preserve.

\begin{table}[h]
\centering
\caption{Failure probability under the limiting distribution $f(\infty)$
(minimum variance $\sigma^2=0.5$) as $k$ increases, uncompressed. The
$2^{-128}$ bound is first violated at $k=30$, so $k=29$ is the ceiling of
the $f(t)$ family at $n=256$, $q=3329$.}
\label{tab:theoretical-ceiling}
\begin{tabular}{cccc}
\toprule
$k$ & Failure probability & pk (bytes) & ct (bytes) \\
\midrule
24 & $2^{-158.7}$ & 9248  & 9600  \\
26 & $2^{-146.2}$ & 10016 & 10368 \\
28 & $2^{-135.5}$ & 10784 & 11136 \\
29 & $2^{-130.6}$ & 11168 & 11520 \\
30 & $2^{-126.1}$ & 11552 & 11904 \\
\bottomrule
\end{tabular}
\end{table}